% ═══════════════════════════════════════════════════════════════════════
\section{Les transformations dans l'enseignement secondaire québécois}
% ═══════════════════════════════════════════════════════════════════════

\begin{frame}{Cette section vous concerne le plus directement}
  \small Trois objectifs : situer les transformations dans le
  programme québécois, traduire la notation du secondaire dans le
  langage matriciel de ce chapitre, et recenser ce que le point de vue
  supérieur apporte concrètement dans une classe.
\end{frame}

\subsection{Où sont-elles dans le programme ?}
\begin{frame}{Table 1.3 --- Parcours dans le curriculum québécois}
  \tiny
  \begin{center}
  \begin{tabular}{p{2.6cm}p{8.6cm}}
    \toprule
    Niveau & Ce qui est en jeu \\
    \midrule
    Primaire, $3^{\text{e}}$ cycle
      & Réflexion, axe de symétrie ; frises et dallages ; vocabulaire
        informel, pliage et miroir. \\
    $1^{\text{re}}$ secondaire
      & Translation, rotation, réflexion ; construction de l'image
        d'une figure ; figures isométriques. \\
    $2^{\text{e}}$ secondaire
      & Homothétie ; figures semblables ; rapport de similitude. \\
    $3^{\text{e}}$ secondaire
      & Conditions minimales d'isométrie et de similitude ; rapports
        $k,k^2,k^3$ pour longueurs, aires, volumes. \\
    $4^{\text{e}}$ secondaire
      & Géométrie analytique : distance, point de partage, droites. \\
    $5^{\text{e}}$ secondaire (SN)
      & Vecteurs : la translation devient enfin \emph{un vecteur}. \\
    \bottomrule
  \end{tabular}
  \end{center}
  \normalsize
  \vspace{4pt}
  \small \textbf{Deux observations :} (1) même ordre que la hiérarchie
  de Klein (fig.~1.3) --- isométries, puis similitudes, jamais l'affine
  général. (2) La synthèse n'a \textbf{jamais lieu} au secondaire :
  transformations et vecteurs restent deux mondes séparés. C'est vous
  qui apporterez ce lien.
\end{frame}

\subsection{Traduire la notation du secondaire}
\begin{frame}{Table 1.4 --- Dictionnaire des notations}
  \small
  \begin{center}
  \begin{tabular}{lll}
    \toprule
    Secondaire & Se lit & Représentation analytique \\
    \midrule
    $t_{(a,b)}$ & translation de $(a,b)$
      & $\col{x}{y} \mapsto \col{x}{y} + \col{a}{b}$ \\[6pt]
    $r_{O,\theta}$ & rotation de centre $O$, angle $\theta$
      & $\col{x}{y} \mapsto r_\theta\col{x}{y}$ \\[6pt]
    $r_{M,\theta}$ & rotation de centre $M$
      & $r_\theta\left(\col{x}{y} - \vec{m}\right) + \vec{m}$ \\[6pt]
    $s_{d}$ & réflexion d'axe $d$
      & $R_\theta\col{x}{y}$ si $d$ passe par $O$ \\[6pt]
    $h_{O,k}$ & homothétie de centre $O$, rapport $k$
      & $k\col{x}{y}$ \\[6pt]
    $h_{M,k}$ & homothétie de centre $M$
      & $k\col{x}{y} + (1-k)\vec{m}$ \\
    \bottomrule
  \end{tabular}
  \end{center}
  \vspace{4pt}
  \small \textbf{Seule ligne à retenir :} $s_d$ (« symétrie ») au
  secondaire = $R_d$ (« réflexion ») dans ce cours.
\end{frame}

\begin{frame}{Pour la classe 1.9 --- Une notation à défendre}
  \begin{Classe}{Une notation à défendre}
    \small On reproche parfois à cette notation d'être lourde. Elle a
    pourtant un mérite considérable : elle traite la transformation
    comme un \textbf{objet} qu'on peut nommer, comparer, composer ---
    plutôt qu'une action qu'on exécute. C'est le saut conceptuel que la
    didactique identifie comme le plus difficile : passer de la
    transformation-\emph{procédé} à la transformation-\emph{objet}.
    Résistez à la tentation de l'abandonner.
  \end{Classe}
\end{frame}

\subsection{Ce que le point de vue matriciel vous donne}
\begin{frame}{Six choses que le point de vue matriciel change (1/2)}
  \small Vous n'enseignerez pas de matrices au secondaire. Mais :
  \begin{enumerate}
    \item \textbf{Une seule idée au lieu de quatre.} Translation,
      rotation, réflexion, homothétie : quatre cas de
      $w\col{x}{y}=A\col{x}{y}+\vec t$. « Existe-t-il d'autres
      transformations ? » --- vous savez que oui, et lesquelles.
    \item \textbf{Pourquoi l'ordre compte.} $r\circ R \neq R\circ r$ :
      un fait d'\textbf{algèbre} (produit de matrices non commutatif),
      pas une bizarrerie géométrique.
    \item \textbf{D'où vient le $k^2$.} Le déterminant --- ou,
      simplement, une aire = produit de deux longueurs.
    \end{enumerate}
\end{frame}

\begin{frame}{Six choses que le point de vue matriciel change (2/2)}
  \begin{enumerate}
    \setcounter{enumi}{3}
    \item \textbf{Pourquoi les conditions minimales fonctionnent.}
      \textsc{cac}, \textsc{aca}, \textsc{ccc}, \textsc{aa}\dots :
      mystérieuses comme axiomes. Le théorème~1.15 en donne la raison
      --- une application affine est déterminée par 3 points, une
      similitude par 2.
    \item \textbf{Le rôle central de la réflexion.} Le corollaire~1.14 :
      la réflexion engendre tout. Justifie de lui donner plus de
      temps, et de manipuler des miroirs.
    \item \textbf{La liste des trois transformations est incomplète.}
      La réflexion glissée existe --- savoir ce qui manque évite une
      fausseté.
  \end{enumerate}
\end{frame}

\subsection{Les difficultés d'apprentissage, et d'où elles viennent}
\begin{frame}{Des obstacles bien documentés (1/2)}
  \small Pas des erreurs de calcul : des \textbf{conceptions
  cohérentes mais fausses} que les élèves --- et souvent les futurs
  enseignants --- se construisent.
  \begin{enumerate}
    \item \textbf{Transformation vue comme un mouvement de la figure.}
      Remède : demander l'image de points isolés, et du quadrillage
      lui-même.
    \item \textbf{Translation confondue avec « bouger sans tourner ».}
      Remède : construire la translation par les seules coordonnées.
    \item \textbf{Centre de rotation toujours placé sur la figure.}
      Échec massif quand le centre est extérieur. Remède : la
      conjugaison (prop.~1.5), « je vais au centre, je tourne, je
      reviens ».
  \end{enumerate}
\end{frame}

\begin{frame}{Des obstacles bien documentés (2/2)}
  \begin{enumerate}
    \setcounter{enumi}{3}
    \item \textbf{Homothétie réduite à un agrandissement.} $k<0$
      n'est pas imaginé. Remède : tracer des rapports négatifs,
      comparer systématiquement longueurs/aires/angles.
    \item \textbf{« Semblable » compris comme « qui se ressemble ».}
      Remède : la définition 1.11 elle-même --- \emph{toutes} les
      distances multipliées par le \emph{même} nombre.
    \item \textbf{Niveau de van Hiele.} Le raisonnement visuel
      (« ça a l'air pareil ») ne produit pas de justification fondée
      sur les propriétés, même avec les bonnes réponses.
  \end{enumerate}
  \begin{Classe}{Un test à faire passer}
    \footnotesize « Rotation de $90°$ autour d'un point $M$ hors du
    triangle $ABC$. $M$ a-t-il une image ? » --- révèle presque à coup
    sûr la conception « mouvement ».
  \end{Classe}
\end{frame}

\subsection{Des activités qui tiennent la route}
\begin{frame}{Cinq activités}
  \begin{enumerate}
    \item \textbf{Deux miroirs.} Angle $\varphi$, compter les images ---
      découverte expérimentale du théorème~1.13.
    \item \textbf{La chasse aux symétries.} Logos, carrelages,
      flocons : compter rotations et réflexions.
    \item \textbf{Le pantographe.} Homothétie manipulable (prop.~1.5) ;
      « où est le centre ? » devient concret.
    \item \textbf{Programmer une transformation.} Écrire la fonction
      qui calcule l'image d'un point --- force à expliciter que la
      transformation agit sur \emph{tous} les points.
    \item \textbf{La géométrie dynamique.} Construire, puis
      \emph{déplacer} un sommet : ce qui reste invariant est ce que la
      transformation conserve.
  \end{enumerate}
\end{frame}

\begin{frame}{Pour la classe 1.11 --- Une mise en garde sur les outils}
  \begin{Classe}{Une mise en garde sur les outils}
    \small La géométrie dynamique et la programmation ne produisent
    pas automatiquement de la compréhension. Un élève peut faire
    glisser un curseur vingt minutes sans rien conclure. Ce qui fait
    apprendre, c'est la \textbf{prédiction} : écrire à l'avance ce que
    l'on croit obtenir, puis exécuter et comparer. Sans cette étape,
    l'outil ne fait que \emph{montrer} ; avec elle, il
    \emph{confronte}.
  \end{Classe}
\end{frame}
